\section{Finite arithmetic, reproduction, and conclusion}\label{sec:verification}

\subsection{Indexing and arithmetic in the certificate program}
Appendix~\ref{app:certificate} lists the complete program used for
Verification~\ref{ver:certificates}; it has no external data input.
Its partition order reads parts from last to first. If \(p_j(u)\)
counts partitions of \(u\) into at most \(j\) parts, the prefix table
counts those whose last padded part is \(<k\), by
\[
 N_j(u,k)=N_j(u,k-1)+p_{j-1}(u-j(k-1)).
\]
The second summand subtracts the last part from all \(j\) positions.
Negative arguments count as zero and \(p_0(0)=1\).
Subtracting two prefix counts ranks the next part after the already
fixed smaller neighbor. A shorter partition padded by zeros has the
same rank. For the tail-restricted table, first group by its exact last
three parts, subtract their largest part from the preceding three, and
use the three-part rank. The grid similarly groups by \(w(m)\) and ranks
the partition corresponding to its first five gaps.

The allocated prefix array has size parameter \(1250\). Its terminal
weight need not supply a complete unrestricted count; all counts used
for multiplicities and grid indices have weight at most \(894\), strictly
below that terminal boundary. The prefix entries are bounded by
\(\binom{1254}{5}<2^{63}\).
The actual largest complete levels used in six, four, and three variables
have respectively \(1229120\), \(261072\), and \(67051\) partitions.
The largest tail-restricted level has \(1263138\) entries, and the grid
has \(6611697\). These ranks, offsets, and predecessor indices fit signed
32-bit integers. Independently of these evaluated counts, coarse bounds
for the three complete levels, the grid, and any tail level are
\[
 \max\left\{\binom{155}{5},\binom{333}{3},\binom{896}{2},
             \binom{154}{5},\binom{21}{3}\binom{272}{2}\right\}<2^{31}.
\]
These count unrestricted weak compositions. For the grid,
\(w(m)\le149\) implies \(\sum_i m_i\le149\); for a tail level there
are at most \(\binom{21}{3}\) ordered tails of size at most \(18\), and
each remaining three-part prefix has weight at most \(270\).

Multiplicity arithmetic is signed 128-bit. The target recurrences check
every accumulation for overflow, and test integrality and nonnegativity
before division. For the source, only \(b\le25\) is needed.
Each intermediate multiplicity is at most \(151^{10}\): it is bounded
by its multiplicity in \((\Sym^b V)^{\otimes j}\), \(j\le6\), and the
successive interlacing partitions have at most
\(151^{0+1+2+3+4}\) choices. In one source Newton step, there are fewer
than \(6\cdot151^4+6\) strip terms, by fixing one coordinate with the
size constraint. Thus every signed accumulation is bounded absolutely by
\[
                   151^{10}(6\cdot151^4+6)<2^{127}.
\]
For \(U\), every width in every actual or high-residue call is at most
\(920\). A chain has at most \(1+2+3+4=10\) width factors, so
\[
                      B_I\le920^{10}<Z=2^{100}.
\]
The minimum initialized at \(Z\) therefore never clips an upper bound,
and summing class weights costs a factor at most \(720\).

For the power bound, \(S\) has
\(\dim\Sym^6\CC^6=\binom{11}{5}=462\) polynomial generators.
Consequently \(\binom{d_0+461}{461}\ge Q(x)\), so the binomial search
for \(z\), before its cap at \(18\), stops by \(461\).
With \(C\le Z\) and \(d_0\le149\), its first crossing is at most
\(150Z\); predivision multiplication costs at most another factor \(610\),
still below \(2^{127}\). After the cap, the other binomials are at most
\(\binom{167}{18}\). Products in the final lower bound are screened
against \(Z\) before multiplication. Capping can only weaken a lower
certificate. These observations cover the unchecked arithmetic as well
as the program's explicit assertions.

OpenMP distributes distinct coefficient destinations within one degree.
All previous degrees are read-only, and each work-sharing loop has its
barrier before the next degree. The mathematical result is therefore
independent of the thread count.

\subsection{Integer bounds for the band program}
It remains to bound the signed 128-bit arithmetic before the
arbitrary-precision division described in Section~\ref{sec:band}.
For a Laurent series truncated to total tail degree at most \(27\),
use the sum of absolute coefficients as its norm.
The norm of \(L_j\) is bounded by
\[
 N=\sum_{n=0}^{27}[z^n]P(z),\qquad
 P(z)=\prod_{h=1}^6(1-z^h)^{-M_h},\qquad
 M_h=(7-h)\binom{h+3}{3}.
\]
Here \((M_1,\ldots,M_6)=(24,50,80,105,112,84)\).
There is a convenient purely elementary estimate
\begin{equation}\label{eq:normbound}
                         N\le4^{27}P(1/4)<2^{79}.
\end{equation}
The first inequality follows from nonnegative coefficients.
For the second, put \(r_h=4^h/2\). Bernoulli's inequality gives
\((1-4^{-h})^{r_h}>1/2\).
Thus the \(h\)-th factor of \(P(1/4)\) is at most
\(2^{\lceil M_h/r_h\rceil}\), with strict total inequality.
The six ceilings are \(12,7,3,1,1,1\), whose sum is \(25\);
combine this with \(4^{27}=2^{54}\).

Each four-coordinate source product has at most \(24\) geometric factors,
each consuming at least one tail unit, so its norm is at most
\[
                         S_0=\binom{51}{24}<2^{51}.
\]
Using \(\|D\|_1\le64\), the quotient \(E_I\) has norm at most
\[
             Q_0=64\binom{21}{6}=3472896<2^{22}.
\]
Indeed dominate all reciprocal cycle factors by \((1-q)^{-6}\) and
sum through degree \(15\). The partial quotient calculations through
degree \(21\) have bound \(64\binom{27}{6}=18944640\), and the
one-coordinate arrays \(A_J\) have coefficients at most
\(\binom{32}{5}=201376\). Their additions fit signed 32-bit integers.
Class-size divisions are exact: for each partial cycle multiset of
size \(r\le6\), its centralizer order divides \(r!\), hence divides \(720\).

The undivided Euler step costs at most \(27N\).
The absolute weight of all target elementary masks is \(720\cdot64\),
and the class-subset weight is
\[
 \sum_{I\vdash6}n_I2^{v(I)}
  =\sum_{\sigma\in\SSS_6}2^{\#\text{cycles}(\sigma)}
  =2\cdot3\cdot4\cdot5\cdot6\cdot7=5040.
\]
The last identity follows inductively by inserting the new largest
element into an existing cycle or making it a new cycle.
Cross multiplication including \(1-q\) has norm at most \(512\);
alternation costs at most \(24\). A bound for every predivision
coefficient, partial sum, and product is therefore
\begin{equation}\label{eq:allbounds}
        512\cdot24\bigl(27\cdot720\cdot64N+5040Q_0S_0\bigr)
                         <2^{115}<2^{127}.
\end{equation}
The source convolutions have nonnegative factors before their final
signs, so this also bounds their intermediate products.
Potentially large products in the code already have a 128-bit operand.
The row divisions and final sums then use arbitrary precision.
Each output counter covers at most six degrees, \(1908\) tails, and
\(448\) second parts per degree and tail. Thus even the coarse bound
\(1908\cdot6\cdot448<2^{31}\) keeps these counters within signed
32-bit range.

\subsection{Recorded computations and reproduction}
The reference output streams contain \(20\) base rows, \(21\) flag rows,
and \(21\) band rows. Tables~\ref{tab:base} and~\ref{tab:band} present
all \(62\) numeric rows, including every tested degree. The reproduction
script checks these streams against fixed hashes and compares their
numeric contents with both manuscript tables.

A full reproduction compiled the two printed sources with
GNU C++~13.3.0 and executed both programs over their entire stated ranges,
with assertions enabled. The first used two OpenMP threads and the
second was serial. Both exited with status zero, and all four output
streams matched the reference streams byte for byte. This execution
reproduced all \(62\) numeric rows in the tables.

The programs retain the complete ranges \(6\le b\le149\).
For the main proof using Corollary~\ref{cor:quadratic-range}, only the
base comparisons \(6\le b\le25\) and the first row of
Table~\ref{tab:band}, for \(26\le b\le29\), are required.
That row contains \(29814\) flags and \(467068\) band coefficients;
the band differences are nonnegative, with \(2114\) equalities.
The later rows give the additional finite checks for the independent
route through Proposition~\ref{prop:large}. Exact target tables at higher
degrees are computed from the recurrences, not assumed from Foulkes'
conjecture, so their use in lower certificates introduces no circularity.

\begin{table}[ht]
\centering
\begin{tabular}{r|r|r|r|r}
\(b\) range&flagged pairs&largest \(t_2\)&band coefficients&zero differences\\ \hline
26--29&29814&25&467068&2114\\
30--35&46805&25&872322&3621\\
36--41&50091&25&1078386&4161\\
42--47&54443&25&1284450&4701\\
48--53&63027&25&1490514&5241\\
54--59&74729&20&1696578&5781\\
60--65&90684&18&1902642&6321\\
66--71&109505&18&2108706&6861\\
72--77&130920&19&2314770&7401\\
78--83&155274&20&2520834&7941\\
84--89&182871&21&2726898&8481\\
90--95&214444&22&2932962&9021\\
96--101&248073&23&3139026&9561\\
102--107&285382&24&3345090&10101\\
108--113&324614&25&3551154&10641\\
114--119&366623&26&3757218&11181\\
120--125&410723&26&3963282&11721\\
126--131&457715&26&4169346&12261\\
132--137&506088&26&4375410&12801\\
138--143&557588&26&4581474&13341\\
144--149&610285&27&4787538&13881\\ \hline
total&4969698&27&57065668&177134
\end{tabular}
\caption{Full retained verification range. The first row is the residual
range required with quadratic stabilization; subsequent rows support the
independent chart route. All flags lie in the checked band, and no
negative difference is found. The program's first row is labeled \(24\), the
multiple of six below the first tested degree; the explicit degree
ranges here avoid that convention.}\label{tab:band}
\end{table}

The supporting package contains the complete program sources, reference
output streams, their hashes and numeric tables, and a reproduction
script. A new full run creates its own compiler and execution metadata.
The commands, run from the paper directory, are
\begin{verbatim}
python3 -B verification/verify_computations.py --check
python3 -B verification/verify_computations.py --run --threads 2
\end{verbatim}
The first checks source and reference-stream identities, all \(62\)
numeric rows, agreement with the manuscript tables, and the numerical
bounds above; it does not execute the mathematical programs.
The second compiles both unchanged sources with
\texttt{-std=c++17 -O3 -UNDEBUG}, enables OpenMP only for the first,
executes them, and requires both complete output streams to match
the reference streams byte for byte. It keeps each run in a new external
directory, with the compiler version, build and execution commands,
exit statuses, timings, hashes, and captured streams. The option
\texttt{--serial} may replace \texttt{--threads 2} to compile and run
the first program without OpenMP.
The two C++ files are also printed in full below, so every finite step
is specified in this paper. The coefficient formulas, indexing
arguments, and arithmetic bounds establish what they check; successful
execution supplies the finite nonnegativity assertions.

\subsection{Completion of the proof}
\begin{proof}[Proof of Theorem~\ref{thm:main}]
Corollary~\ref{cor:quadratic-range} settles \(b\ge30\).
Verification~\ref{ver:certificates} settles \(6\le b\le25\).
Proceed by increasing \(b\) from \(26\) through \(29\).
If \(d_6\ge6\), Lemma~\ref{lem:rectangles} reduces the comparison to
degree \(b-6\ge20\). Otherwise \(d=r+6m\) lies in the exhaustive
classification of Section~\ref{sec:certificates}. Every pair is either
cleared there by a proved certificate or the chart, or flagged.
The flags have \(t_2\le27\) and are covered by
Verification~\ref{ver:band}. Hence \(Q(\lambda)\ge F(\lambda)\) for
every source partition and every \(b\ge6\).
The source has no constituents of length greater than six;
specialization at extra variables zero transfers the comparisons to
every finite dimension. Complete reducibility gives the required
equivariant injection.
\end{proof}

\begin{remark}[The independent full-range route]
The same finite argument, with no change of predicate or arithmetic,
continues from \(b=30\) through \(b=149\), as recorded in the full
tables. Combining it with Proposition~\ref{prop:large} for \(b\ge150\)
proves the theorem without using the quadratic companion.
The chart and certificate methods are therefore retained as a complete
independent proof, while Corollary~\ref{cor:quadratic-range} gives the
shorter large-degree route.
\end{remark}
